In this chapter, I address the first facet of multifaceted geospatial data: its \emph{multivariate structure}.
%
Geospatial data rarely arrives as a single quantity per location; instead, many interacting variables coexist over the same geography, yet the map graph schema that makes such data legible reserves spatial position for location itself, leaving little room to encode additional variables without breaking the viewer's expectations.
%
Compounding this, real-world geospatial data is discretized in irregular and inconsistent ways, raising statistical hazards such as the Modifiable Areal Unit Problem.
%
To confront these challenges, I present \emph{HexTiles}, an attempt at a domain-agnostic hexagonal-tiling encoding that visualizes multivariate geospatial data and their spatial interactions simultaneously while remaining robust to the differing discretizations geospatial data arrives in.
%
I evaluate HexTiles through a controlled user study and a domain-expert review with ecologists and hydrologists, and the remainder of this chapter develops the motivation, design, and assessment of the approach in detail.
\section{Introduction}\label{sec:intro}
Multivariate visualization, i.e. visualizing multiple variables  or information from multiple information channels together, is a recurring problem that the visualization community has sought to solve \cite{liu2014survey}. 
%
% In fact, the study in dimensionality reduction and visualization methods that leverage these techniques arise from this exact challenge \cite{espadoto2019toward}. 
% Given the ever-increasing data size and dimensions, this continues to be a challenge.
%    
This is no different for geospatial data \cite{andrienko_geovisualization_2008, guo_multivariate_2005, strode2020bivariate, saha_visualizing_2022}. 
%
The challenge of visualizing and communicating insights from multivariate geospatial information data has presented, and still presents, a unique challenge in visualization and cartography~\cite{bertin1983semiology, Tyner_2014}.
% , as the range of geospatial visualizations can be limited due to the rigid graph schema of geospatial maps.
%
Pinker, in their seminal work, writes that ``people create schemas for specific types of graphs using a general graph schema, embodying their knowledge of what graphs are for and how they are interpreted in general \cite{pinker1990a}.''
%
For example, in a bar graph, the dependent variables and the independent variables tend to be on the Y-axis and X-axis, respectively.
%
In interpreting a bar graph, we \textit{match} this prior experience and expectation to the visuals. 
%
This \textit{match}, or lack thereof, is a key component in how we reason about and make decisions when using visualizations \cite{padilla_decision_2018}.
%
In contrast to other graph schemas, maps, where most geospatial data is visualized, have a comparatively stricter set of expectations that most likely arises from everyday experience with maps including the cardinal directions~\cite{hografer2020state}.
 
For geospatial visualization designers, this is a two-edged sword. 
%
This stricter schema implies that if the visualization appropriately \textit{matches} this graph schema, the viewer can quickly start the \textit{interrogation} process where the visual data is converted to answer a particular \textit{conceptual question}.
%
On the other hand, the strict nature of the schema can also partially limit the design space of geospatial visualizations, as the spatial position of the visualization interface is often reserved by the geographic position.
%
% On the other hand, the nature also means that departures from this map graph schema, will render the visualization undecipherable as viewers are essentially required to fight the preexisting expectations.
%
This may be a part of the reason why methods and techniques in geospatial visualization have largely unchanged, especially for visualization used for the general public.
%%
%The widespread adoption of the choropleth map is exemplified by a study that examined state COVID dashboards for all 50 US states that found 94\% of the dashboards used a choropleth map \cite{adams_normalizing_2023}.
%
% While the choropleth map can be a powerful visualization tool, it has a number of limitations including its challenge with encoding more than two variables\cite{Tyner_2014}, which we solve with our proposed HexTiles design.

Furthermore, multivariate geospatial visualizations must also properly contend with data with varying spatial discretizations.
% \textbf{Spatial discretization mismatches}
% \note{add points too}
%
% Here, spatial discretization refers to how geospatial data \note{wording}
%
Unlike scientific data, where a simulation code may solve a series of PDEs on a regular grid, geospatial data is often discretized in irregular patterns, following geographical features (lakes, mountains, etc.) or based on other geographical groupings like counties. 
%
% Beyond polygon-based data, data can also be formatted as a set of points, where each coordinate records some geospatial event, as an example.
%
As such, it is commonplace for geospatial data visualizations to work with data using varying spatial discretizations.
%
However, differing spatial discretizations also leads to issues like the Modifiable Areal Unit Problem (MAUP) \cite{fotheringham1991modifiable} which we discuss in \cref{sec:methods} and our encoding approach to mitigate its effects.
%%
%This, thus excludes methods like the choropleth map as a general tool for multivariate geospatial visualization.
%%
%We further describe limitations of colors in maps in \cref{sec:limitations} and how we solve this problem with spatial binning in \cref{sec:methods}.

In order to address these challenges, we propose HexTiles, a domain-agnostic hexagonal-tiling based visual encoding design for multivariate geospatial data.
%
In short, the design of HexTiles is motivated by the need to (1) faithfully visualize geospatial data of varying discretizations and qualities simultaneously, 
and (2) to highlight the interactions and spatial relationships between variables across spatial regions.
%
We assess the efficacy of HexTiles via a user study and a domain expert review with ecologists and hydrologists experienced in designing geospatial visualizations.

% As the scale of data grows, the need for accurate decision making based on these data also grows together.

% The paper is organized as follows.
% %
% We first cover the related work in~\cref{sec:relworks}, describe the methods we apply in~\cref{sec:methods}, the resulting visual encoding in~\cref{sec:design}, the case studies in~\cref{sec:cases} and finally the expert feedback in~\cref{sec:feedback}.




%\begin{figure}[htb]
%  \centering
%  % the following command controls the width of the embedded PS file
%  % (relative to the width of the current column)
%  \includegraphics[width=0.8\linewidth]{chart}
%  % replacing the above command with the one below will explicitly set
%  % the bounding box of the PS figure to the rectangle (xl,yl),(xh,yh).
%  % It will also prevent LaTeX from reading the PS file to determine
%  % the bounding box (i.e., it will speed up the compilation process)
%  % \includegraphics[width=.95\linewidth, bb=39 696 126 756]{sampleFig}
%  %
%  %\parbox[t]{.9\columnwidth}{\relax
%  %         For all figures please keep in mind that you \textbf{must not}
%  %         use images with transparent background! 
%  %         }
%  %%
%  \caption{A choropleth map of the United States that encodes data in a sequential discrete colormap. From: \url{https://observablehq.com/@d3/choropleth}}
%  \label{fig:chloro}
%\end{figure}

\section{Related Work}\label{sec:relworks}
Our work is related to others exploring multivariate geospatial visualization and its methods. 
%
% Here, we discuss the body of past work in this area, as well as the limitations of prior approaches that we identify.
Here, we discuss the body of past work in this area in visualization as well as cartography.

% \subsection{Multivariate Geospatial Visualizations}

%Similarly, when we read geospatial maps, viewers have a set of prior expectations of what a map should look like - weexpectation of cardinal directions to orient a map or the expectation on inherent ordering of spatial descriptors (e.g. latitude and latitude) are examples.
%%
%As Hogr{\"a}fer et. al. write maps must support the view's kl 




%
%\note{Expand this} \note{The map is often seen, therefore, as a metaphor for the presentation and representation of all types of data, spatial and non-spatial. - Fairbairn 2001} Geospatial data, especially to the general public, is only represented on maps. 
%%
%Of course, minor distortions of maps have proven to be effective in some cases, \note{Metro line maps?\cite{nesbitt2004getting}} but we find that geospatial data give little leeway for this reason. 
%%
%Visualizing this data via other representations often contradicts the graph schema - we do not expect geospatial data to be represented without a map.



% \subsubsection{Unique challenges in geospatial visualizations}

Previous authors \cite{andrienko_geovisualization_2008, guo_multivariate_2005,mcnabb_multivariate_2019, kim_bristle_2013} have identified a clear need for finding visualizations that facilitate the examination of two or more spatially bound data classes \cite{guo_multivariate_2005}.
%
As Javed and Elmqvist \cite{javed_exploring_2012} describe in their survey of composite visualizations, four approaches exist for combining visual representations: juxtaposition, superimposition, overloading, and nesting. 
%
Similar taxonomies to organize geospatial visualization have been proposed by previous works. 
%
Many works \cite{andrienko2010space, mota_comparison_2022, mayr_once_2018} have compared the efficacy of the different approaches of visualizing multivariate information across varying applications \cite{pena-araya_comparison_2020, griffin2006comparison}.
%
For geospatial visualizations, however, juxtaposition and superimposition are the common techniques for composing information channels.

Cartography has also extensively studied geospatial visualizations, most notably by Bertin \cite{bertin1983semiology} who in his seminal 1963 text proposed a set of visual variables that can be used to encode data on maps. 
% \fix{Need more cartography works here}
%
The ubiquitous choropleth map, for example, is a direct application of Bertin's visual variables where color is used to encode a variable of interest.
%%
%Color is a highly used visual variable in geospatial visualizations, used both for categorical and continuous data.
%
Colors can also be used to encode multiple variables, as in bivariate colormaps, where two variables are encoded in the hue and saturation of a color \cite{Tyner_2014}.
%
%Bivariate colormaps and its efficacy have also been studied, including guidelines for their design, by many authors
Bivariate colormaps and its efficacy, including guidelines for their design, have also been studied by many authors~\cite{eyton1984complementary, olson1981spectrally, wainer_empirical_1980, strode2020bivariate, schloss_mapping_2019}. 


%We also want to be precise about what multivariate geospatial visualization encompasses.
%%
%Typically, the term \textit{multivariate} data is characterized by the number of keys versus values \cite{Munzner_2015}.
%%
%However, multivariate geospatial data take on many forms. 
%%
%% An election data set, for example, can describe each county with the shares of votes for some party A, its population, median income, etc.
%%
%Uncertainty visualization for geospatial data is one instance of a multivariate data and remains a highly active field in which many approaches have been proposed \cite{maceachren2005visualizing,preston_uncertainty-aware_2019,mckenzie2016assessing}.
%%
%% For many data sets, data values are measured by devices (e.g. by a satellite); hence, some degree of uncertainty may be associated with some data values.
%%
%In addition to these cases where the data can be described in a tabular format, time-varying data can also considered a multivariate data with its inclusion of the temporal axis.
%%
%Geospatial-temporal visualization and effective strategies are also studied \cite{andrienko2010space,mota_comparison_2022,ondov_coronaviz_2022}; however, we exclude this from this work and leave that as further work. 
% \note{ Although, we are concerned with multivariate (many different information), time can also be considered an information channel.  }

%\begin{figure}[htb]
%% \begin{figure}[]
%  \centering
%  % the following command controls the width of the embedded PS file
%  % (relative to the width of the current column)
%  \includegraphics[width=0.85\linewidth]{JuxVsSup}
%  % replacing the above command with the one below will explicitly set
%  % the bounding box of the PS figure to the rectangle (xl,yl),(xh,yh).
%  % It will also prevent LaTeX from reading the PS file to determine
%  % the bounding box (i.e., it will speed up the compilation process)
%  % \includegraphics[width=.95\linewidth, bb=39 696 126 756]{sampleFig}
%  %
%  %\parbox[t]{.9\columnwidth}{\relax
%  %         For all figures please keep in mind that you \textbf{must not}
%  %         use images with transparent background! 
%  %         }
%  %%
%  \caption{Comparing the two general strategies for multivariate geospatial visualizations: Juxtaposition (left) and Superimposition (right). Each variable is represented by a different color.}
%  \label{fig:jux_vs_sup}
%\end{figure}

% \subsection{General approaches}
Generally, there are two major approaches for multivariate geospatial visualizations: \textbf{juxtaposition} and \textbf{superimposition}, which we describe in order in \cref{subsec:jux,subsec:super}.
%
However, we note that, thanks to the large body of past work in multivariate geospatial visualizations, a number of works have explored other visualization designs beyond juxtaposition and superimposition.
%
Methods that leverage textures \cite{maceachren1998visualizing}, density via stippling \cite{gortler_stippling_2019}, and density of lines or \textit{bristles} \cite{kim_bristle_2013} are among the other many methods proposed.
% \fix{Some other methods, but we want to focus on Juxtapotion and Superimposition here as general strategies}

%These include methods like SOM (self-organizing maps) \cite{kohonen_self-organizing_1990}, and the many visualization tools that leverage this technique \cite{takatsuka2001application,akande2017geospatial}.
%%
%However, we find that these can be inaccessible tools for the general public and are not readily decipherable, and thus limit the scope here to juxtaposition and superimposition.

\subsection{Juxtaposition}
\label{subsec:jux}
Javed and Elmqvist highlight that juxtaposition is the most common and flexible approach for combining visual visualizations \cite{javed_exploring_2012}.
%
The small multiples approach or linked views approach, which leverage more interaction for users to better reason about the juxtaposed geospatial visualizations, are frequently employed \cite{deng_visualizing_2023, yang2022epimob}.
%
In many cases, these designs envision to link map-based geospatial and temporal data with other common data encodings like line charts, etc., but the key thread remains to encourage the users to reason about the presented data via implicitly or explicitly comparing the data shown in each visualization.
%
% For the multivariate geospatial case, a set of small multiples plots each displaying an unique variable is one such example.
%
% \note{fix}
However, despite its simplicity, juxtaposition has been showed to be ineffective at times \cite{ramachandran2009visualizing, nusrat2017cartogram}.
%
Increasing the number of views can render the contextualization of a larger number of views difficult, since viewers must switch back and forth between views to compare the data.
% as constantly having to switch between views increases the cognitive load put on a user.
%
This is a major limitation of this approach.

\subsection{Superimposition}
\label{subsec:super}
Superimposition refers to displaying multiple data channels in the same visualization instance at once, and can take a number of forms, including \textbf{layering} and \textbf{glyphing}.
% \note{need to define and cite the definition of layering and glyphing}

\subsubsection{Layering}
\textbf{Layering} and its use is heavily studied and explored in many different domains \cite{allan_decal_2017, lobo2017mapmosaic, heer2009sizing}.
%
An often used analogy is the creation of a painting, when an artist composes several consecutive layers of paint from basic background elements to details.
%
Each individual layer can encode one or more variables, and the final composition can be used to communicate the multivariate geospatial data.
%
However, the principle challenge here is to reduce and limit visual clutter and occlusion as noted by several works \cite{mota_comparison_2022, mayr_once_2018}.
%
Careful considerations are required in the visual encoding and usage of color, etc. to ensure that each layer is decipherable and meaningful. 

The simplest method of layering in the context of geospatial data is to combine two or more conventional encodings on top of one another on one map.
%
For example, one can combine a choropleth map with a proportional symbol map, where the size of the symbol encodes another variable to encode two variables at once. 
%
Similar designs but leveraging a pie chart in place of a proportional symbol map along with a choropleth map have been used as well \cite{Tyner_2014, vallandingham}.
%.
These designs that enable users to visually separate the variables are referred to as extrinsic methods as opposed to intrinsic methods like bivariate colormaps.
%
Though both are effective for certain purposes, recent works in cartography have explored the efficacy of the two designs.
%
In fact, for two variables, Stachoň et. al. and previous works demonstrated that extrinsic methods, e.g. combining two variables via layering two separate encodings, can outperform intrinsic methods like bivariate colormaps \cite{stachon_comparison_2023, sasinka_comparison_2021}.
%
The proposed HexTiles design thus also adopts an extrinsic approach to combine multiple variables, rather than envisioning a more visually continuous method like a bivariate or a trivariate colormap.
% However, most use cases for the layering approaches consider two variables as the visual clutter often increases with every additional variable considered.
% For geospatial applications, layering is the most used and generally most intuitive to the general public via its use in public-facing tools like Google Maps and commonly used libaries like \texttt{ggplot}\cite{kahle2013ggmap}.

\subsubsection{Glyphs}
In addition to layering, \textbf{glyphs} are another fundamental technique to encode multivariate data, especially in geospatial contexts \cite{mcnabb_multivariate_2019,ondov_coronaviz_2022}.
%
With the flexibility in mapping choices their design space offers, glyphs have become a highly popular way to visualize multivariate data. 
%
As such, significant research has been dedicated to describing glyphs in the context of visualization. 
%
Among the first, Chen et. al. \cite{chen2008multivariate} details their theoretical principles, such as those pertaining to visual mapping, and techniques for their layout, while addressing important concerns such as the bias certain encodings may introduce.
%
%Ropinski et al. \cite{ropinski2008taxonomy, ropinski_survey_2011} present two surveys that cover the use of glyphs for medical visualizations, in which they focus on the perception of their channels. 
%%
%They differentiate between variables that quickly draw out attention (pre-attentive) and those that require focused attention (attentive). 
%%
%Expanding on these foundations, Borgo et al. \cite{borgo2013glyph} give a holistic discussion of glyphs. 
%%
%Their seminal work offers a rich discussion of a broad, however, highly relevant set of related topics, such as semiotics, or perceptual studies, based on which they give design guidelines for their creation.
%
%Zhang and Pazner \cite{zhang2004icon} were among the first authors to use multivariate glyphs on maps. 
%%
%Their system allows for the creation of glyphs with juxtaposing visual elements. Within our design space, we create glyphs within a regular grid that we base on the findings mentioned above.

Glyphs have also been studied extensively in cartography, and in fact predates the study of glyphs from the current field of data visualization \cite{anderson1957semigraphical,elmer2013symbol}, and have been showed to be effective at communicating multivariate data to audiences with varying levels of expertise and familiarity with the data \cite{strode2020bivariate}.
\begin{figure}
    \begin{center}
        
\includegraphics[width=0.5\linewidth]{4dsquare}
    \end{center}
    \caption{4D square-glyph example \cite{strode_exploratory_2020}.}
    \label{fig:squareglyph}
    \vspace{-0.1in}
\end{figure}

\begin{figure*}[h!]
    \centering
    \begin{subfigure}[b]{0.47\textwidth}
        \centering
        \includegraphics[width=0.7\linewidth]{square-glyph-original}
        \caption{3D square-glyphs used to encode a road network and its three associated attributes in a map of the city of Zurich. From \cite{bleisch2018exploring}. Encoded are, from top-right clockwise: Greenery, Transport Accessibility, and Shop Accessibility.}
        \label{subfig:square_glyph_bleisch}
    \end{subfigure}
    \hspace{2em}%
    \begin{subfigure}[b]{0.47\textwidth}
        \centering
        \includegraphics[width=0.7\linewidth]{square-glyphs}
        \caption{4D square-glyphs used to encode four variables in one map over the Tallahassee area in Florida. From \cite{strode_exploratory_2020}. Encoded are, from top-right clockwise: Household Composition/Diability, Housing/Transportation, Ethnicity/ Language, and Socioeconomic Status.}
        \label{subfig:square_glyph_strode}
    \end{subfigure}
    \caption{Examples of real-life use cases for square-glyphs to encode multivariate geospatial data.}\label{fig:squareglyphexample}
    \vspace{-0.15in}

\end{figure*}



Among the various designs, Square-glyphs as seen in \cref{fig:squareglyph,fig:squareglyphexample} in particular have been used in past works to handle more than two variables \cite{strode_exploratory_2020, bleisch2018exploring}, and was also recently studied by Müller et. al., comparing efficacy of 2D vs. 4D square glyphs \cite{muller_square_glyphs_2023}.
%
Simply put, the square-glyph is a design that employs two visual channels, color and area, to encode up to four variables in one glyph.
%
When used, the underlying geospatial space is divided into a square grids, and whenever applicable, a square-glyph is placed in a grid cell along with each of the up to four variables encoded in the area of the squares as seen in ~\cref{fig:squareglyph}.
%
At the same time, color is used to encode the data dimension so that same color addresses the same variable across all square-glyphs. 
%
Square-glyphs along with radar, or snowflake glyphs, are among the more commonly used glyphs for multivariate geospatial visualizations with more than two variables.

Another class of visualization methods that handle multivariate data are cartograms, which are also known to be effective in communicating geospatial data~\cite{nusrat2017cartogram}.
%
Various flavors of cartograms have been proposed over time, and we refer the reader to a comprehensive review of the past work in this area by Nusrat and Kobourov~\cite{nusrat_state_2016}.
%
However, as Nusrat and Kobourov call attention to, cartograms are also limited in that they inherently distort the base map, despite the strong visual impact they can have for many use cases.
%
Furthermore, they also recognize that multivariate cartograms is a largely unexplored area with further work needed to establish best practices for their design and use.
%
In addition to the visual clutter issue that layering faces, glyphs can also face the issue of placement \cite{mcnabb_multivariate_2019, tsorlini_designing_2017}.
%
McNabb and Laremee offer a comprehensive review of the past work in this area in \cite{mcnabb_multivariate_2019}.
%
However, we still consider this to be an open research question, and instead consider a spatial binning approach that trivializes and simplifies this problem as described in \cref{sec:methods}.

Owing to the documented success of glyphing from past work in visualization and cartography, with HexTiles we adopt a glyph-based approach and develop a superimposition-based visualization design for multivariate geospatial data.

% Directly plotting such multivariate data, where data points may share a spatial location, introduces occlusion as a function of the size of the data set. Some works have examined optimization strategies for
% this\cite{mcnabb_multivariate_2019}. 

% \subsection{Other approaches}
%As with all visualization design, the audience and application domain is a key consideration.
%%
%Naturally, many different disciplines produce different type of data with different needs and a general purpose tool should be able to support a wide range of use cases.
%%

%As previously mentioned, multivariate data need not strictly be data with multiple data variable. 
%%
%For example, uncertainty geospatial visualization for disaster management is an use case with possible life-threatening consequences if poorly executed \cite{gomov_interactive_2019, preston_uncertainty-aware_2019, ruginski_non-expert_2016}.
%%
%Here the uncertainty of a measurement, say of the measured strength of fire at some location, is another set of information that is worth communicating especially in the context of wildfire management.
%% 
%\note{need work}
%Proper handling and encoding of the various data types is critical and these decisions are outlined in \cref{sec:design}.
%% An addition of a temporal axis is another dimension that needs consideration. Animation or a small multiples plots is the current method to solve visualization for this additional dimension, but they both have issues. In short, designing a visual abstraction that effectively communicates multivariate geospatial data on one static interface is an open research question.

%% \note{Motivate the focus for our encoding to be for the general public. For experts, other methods exist}
%We also reiterate that our visual encodings are targeted towards the general public, and those with not necessarily the background knowledge to interpret complex visual encoding.
%%
%In other words, we want to ensure that we do not depart from known visualization styles or graph schemas \cite{padilla_decision_2018}.
%


% \note{need more here, I think}

% \subsubsection{Common approaches in use}

% \textbf{Layering} typically refers to the way data mappings are superimposed into one coherent visualization. 

% \note{Layering in the context of geospatial}


% \note{
%     occlusion: Andrienko et al. directly embedded ThemeRiver plots onto geographical maps, but this method led to severe clutter and occlusion of other map information
% }



% \textbf{VA type works}
% \cite{deng_visualizing_2023}

% \subsection{Unique challenges in geospatial visualizations}

%\textbf{The loss of two dimensions}
%Barring immersive visualization methods that can take advantage of a third dimension, visualization interfaces are always two dimensional, whether it be on print or screen. 
%%
%% This is another challenge that multivariate geospatial visualizations face.
%%
%Since latitude and longitude are data axis that cannot be reordered, a geospatial map inherently occupies the two dimensions of any interface \cite{franconeri2021science}. 
%%
%This contrasts with many other data sources, where the data axis can be appropriately omitted or reordered, while still maintaining key information.
%%
%As such, almost all visualizations from Google Maps to ESRI \note{cite ESRI} rely on a map as a base for geospatial visualizations.
%%
%This can limit many visualization designs since designers are forced into this geospatial map which have strict requirements to ensure interpretability.
%%
%Now, with the widespread adoption of interactive devices like smartphones, visualization design can now leverage 3D graphics; however, the use of the third dimension needs be done with caution for geospatial maps due to occlusion, but more importantly, because any geospatial map that uses the z-dimension (i.e. encode some data in height) competes with prior assumptions that height on maps is used to encode elevation.



% \subsection{On the limitations on color and choropleth maps}
%\section{On color and choropleth maps}
%\label{sec:limitations}
%When considering geospatial visualizations, perhaps the most common channel that visualization designers use is color saturation.
%%
%The commonly seen univariate choropleth map is a prime example (\cref{fig:chloro}).
%%
%The widely used and adopted \texttt{ColorBrewer} tool which recommends color schemes for various data types for geospatial visualization \cite{harrower2003colorbrewer}, is also testament to its popularity.
%%
%Now, as mentioned prior, there is some evidence that extrinsic methods for multivariate visualization can outperform intrinsic methods like bivariate colormaps \cite{stachon_comparison_2023, sasinka_comparison_2021}. 
%%
%Thus, in designing HexTiles, as we considered a multivariate visualization, we made

%The effect of color in geospatial visualizations and colormaps are highly active research areas.
%%
%Works on the efficacy of bivariate colormaps \cite{wainer_empirical_1980}, effective strategies for bivariate colormaps \cite{strode2020bivariate}, effect of color saturation \cite{scholz2014uncertainty}, and the effect of biases in the understanding of color have been studied \cite{schloss_mapping_2019}.

%However, while color hue is known to be effective for categorical data \cite{Munzner_2015,franconeri2021science}, color saturation, for example, is not known to be particularly effective for ordered attributes.
%%
%A natural question to ask is: Why is color saturation so often used for ordered attributes despite its known limitations?

%Put simply, we hypothesize this is primarily because other channels like position and length are not conducive for geospatial maps.
%%
%Due to the inherent rigid structure of geospatial maps, as alluded to in~\cref{sec:intro}, and since the longitude and latitude already takes up the two dimensions available on most mediums barring immersive visualization applications, adding additional visual channels likely complicated an already potentially complicated visualization.
%%
%In addition to this, another possibility is that the general public is too used to interpreting data from color on maps, whether that be a map of your city, a hiking trail map, or an atlas --- color and maps are so strongly linked that it makes the departure of this trend difficult.

%%
%%Cartograms attempt to encode information in 2D area, which is known to be more effective than color saturation for ordered data\cite{franconeri2021science}; however, cartograms are also limited in that they inherently distort the base map and the geospatial information therein.
%%%
%%While this tradeoff is certainly favorable for many cases, we find t
%%%
%%The study of color theory and its implications for data visualization is well documented.

%This is not to say that we advocate for \textit{replacing} color altogether. 
%%
%We consider that the use of color saturation for geospatial maps, comes about as a natural consequence with the limitation of the medium, and can be useful in displaying and communicating general structure of data.
%%
%With HexTiles, in addition to using color, we also encode data and information in icons.
%%
%This allows us to reduce the reliance on colors as well as take advantage of several useful properties of icons as means to represent data as we describe in \cref{sec:methods}.
% In particular, we encode data in the number of icons to ease the sense-making process for colors.

% \nocite{*}

\section{HexTiles}\label{sec:methods}
The design of HexTiles is motivated by the need to (1) faithfully visualize geospatial data of varying spatial discretizations and (2) ease the process of assessing interactions between variables across spatial regions.
%
To this end, taking inspiration from past works in superimposition-based multivariate geospatial visualizations like Square-glyphs, we specifically enumerate the design considerations (DCs) for HexTiles as follows:
\begin{tightItemize}
    \item \textbf{DC1:} Handle geospatial data of varying spatial discretizations simultaneously.
    \item \textbf{DC2:} Allow for assessment of interactions between variables across spatial regions.
    \item \textbf{DC3:} Enable visualizations that allow designers to steer sensemaking and prioritize certain information over others.
    \item \textbf{DC4:} Enable a more faithful visualization of geospatial variables that take into account the error that is introduced when data is spatially aggregated and thus to alleviate the modifiable areal unit problem (MAUP).
\end{tightItemize}
\subsection{Design}
%
To address these challenges, we highlight the three major design decisions for HexTiles and which \textbf{DC} they address: (1) spatial binning [\textbf{DC1}], (2) semantic icons [\textbf{DC2, DC3}], and (3) confidence encoding [\textbf{DC4}].


\subsubsection{Spatial Binning}
\textbf{Spatial binning} is a widely used form of visual abstraction in which data is aggregated within polygonal regions. 
%
This is required as geospatial data is often irregularly discretized.
%
While this approach sacrifices a viewer’s ability to examine individual data points, Cleveland \cite{cleveland1985graphical} highlights the ability of these bins to convey their underlying spatial distributions. 
%
To form bins, both irregular and regular polygons can be used. 
%
% While irregular binning is commonly found when data relates to, for example, political boundaries, Battersby et al. \cite{battersby2017shapes} argue the perceptual advantages of regular bins for non-anthropogenic data.
%
Such regular tessellations, which can only be achieved using either equilateral triangles, squares, or hexagons, produce a smoother gradient and allow for a direct comparison between tiles. 
%
As triangles alternate their orientation within a grid, squares and hexagons are predominantly used for spatial binning on maps.
%
Carr \cite{carr1990looking} explored the use of both shapes in binned data plots for visualizing large datasets.
%
He and his colleagues \cite{carr_hexagon_1992} later compared square and hexagonal tessellations for univariate and bivariate mosaic maps. 
%
% Them, along with several more recent publications \cite{birch2007rectangular, shelton2014mapping} stress the following three major advantages of using hexagonal over square tessellations. 
Them, along with several more recent publications \cite{birch2007rectangular, shelton2014mapping}, advocate for the use of hexagonal tessellations, primarily citing the higher distraction that square tessellations introduce due to the horizontal and vertical lines.
%
% In addition, we consider hexagonal tessellation to be accessible to many audiences, as they are commonly used in board games and other mediums as well.
%stress the following three major advantages of using hexagonal over square tessellations. 
%First, square grids introduce horizontal and vertical lines, which human perception has an innate preference for. 
%%
%This preference may elicit a perceptual bias and may obscure the identification of spatial patterns \cite{carr_hexagon_1992}. 
%%
%Second, Scott \cite{scott1985averaged} found that hexagons have higher representational accuracy, and third, better connective properties–surrounding tiles in a tessellation are equidistant from their center—are all properties attributed to the recent surge in publications using hexagonal tiling \cite{birch2007rectangular, shelton2014mapping}. 
%%
%The benefit of spatial binning is two-fold: (1) by introducing a regular discretization we can solve the spatial discretization mismatch issue from~\cref{sec:intro}, and (2) ease the comparison process between two regions as the unit of comparison becomes the HexTiles, in the case of our work.
%
%As scientific results across domains feature data with a spatial connection, which is increasingly communicated through geovisualizations, more fields have started exploring the use of geographic space for their field, in what has become known as the spatial turn \cite{skupin2009discrete}. 
%%
%Birch et al. \cite{birch2007rectangular} surveyed the use of tessellations in ecology and support the notion of a trend towards favoring hexagonal tessellations. 
%%
%In further support thereof, other examples can be found across non-visualization domains \cite{shelton2014mapping, shelton2015social}. 
%%
%The question of how to best utilize the available cells, however, is a topic of ongoing research. Heimerl et al. \cite{heimerl2018visual} have recently explored the use of encodings and their combination for hexagonal scatterplots.
%%
%They compiled a comprehensive list of tasks commonly found for analyzing these plots and consider designs appropriate for the different analysis questions. 
%
This all gives strong reason to base our visualization on hexagonal tiling.

\subsubsection{Icons}
\textbf{Icons} are a common form of visual abstraction in which data is represented by a symbol, often taking advantage of metaphors  \cite{nocke2005icon, presnov_pacemod_2023}.
%
Our main motivation of using icons to represent data is (1) to leverage the \textit{pre-attentive} processing of visual information, much like what Chernoff attempts via data encoding in facial features \cite{chernoff1973use}, and (2) to reduce reliance on colormaps to encode data.
%
%While a number of considerations must be made when designing icons, one key consideration is to ensure that icons are pre-attentively \textit{distinct} from each other, as pointed out by Ware \cite{ware2019information}.
%%
%% As Ware points out, in designing symbols for visualizations, it is important that symbols are .
%%
%This is our view on why bivariate or multivariate colormaps can be challenging to design well.

 % \begin{wrapfigure}{r}{0.5\linewidth}
 %   \begin{center}
 %     \includegraphics[width=0.5\linewidth]{bivariate}
 %   \end{center}
 %   \caption{A bivariate colormap}
 %   \label{fig:bivariate}
% \end{wrapfigure}

\begin{figure}[h!]
    % \small
    \centering
    \begin{subfigure}[b]{0.4\linewidth}
        \centering
        \includegraphics[width=\linewidth]{IconExample}
    \end{subfigure}
    \begin{subfigure}[b]{0.4\linewidth}
        \centering
        \includegraphics[width=\linewidth]{ColormapExample}
    \end{subfigure}
    \caption{The same CalSim3 unmet water demand data (details in~\cref{sec:cases}) over a part of California's central valley; \textit{Left:} icons, \textit{Right:} continous red colormap. Each icon represents around 0.15 Acre-feet / Acre of unmet water demand within the corresponding region. As seen in the number of icons and the colormap, unmet water demand increases from the HexTile region on the left to the right.}%missing
    \label{fig:icon_example}
\end{figure}
%% \begin{figure}[]
%  \centering
%  % the following command controls the width of the embedded PS file
%  % (relative to the width of the current column)
%  \includegraphics[width=0.3\linewidth]{bivariate}
%  % replacing the above command with the one below will explicitly set
%  % the bounding box of the PS figure to the rectangle (xl,yl),(xh,yh).
%  % \includegraphics[width=.95\linewidth, bb=39 696 126 756]{sampleFig}
%  %
%  %\parbox[t]{.9\columnwidth}{\relax
%  %         For all figures please keep in mind that you \textbf{must not}
%  %         use images with transparent background! 
%  %         }
%  %%
%  \label{fig:bivariate} 
%  \caption{Bivariate colormap}
%\end{figure}

An example of using the number of icons to represent univariate unmet water demand data (details in \cref{sec:cases}) is shown in \cref{fig:icon_example}.
%
In order to ease interpretation, we use a red water droplet and encode the amount of unmet water demand with its count.
%
With reference to \cref{fig:icon_example}, we highlight some qualities of icons as well as some trade-offs.

First, icons can be expressive, since we can almost immediately detect if there are one, two, three or four objects in a group - an ability that we possess from early on in our development \cite{dehaene2011number,ware2019information}.
%
Second, comparing two regions is made easier since we can compare the number of icons in each region.
%
We argue that similar operations can be more difficult with colors.
%
Lastly, and importantly in the context of HexTiles, by using icons we can reduce the burden of color encoding, making way for other variables to be simultaneously encoded in color. 

%In many cases, interpreting data encoded with bivariate colormaps become an \textit{attentive} endeavour, requiring viewers to mentally map back and forth between the color and the data.
%% 
%Even though, some colors have meaningful association to sentiments or other concepts, immediately mapping color from a bivariate colormap to meaning is difficult.
%%
%In contrast, with HexTiles, in addition to color, we also employ icons to represent data.

%For continuous data, this is equivalent to mapping data to an discrete colormap, where instead the we replace the color with a group of icons, where the number of icons is dependent on the value to be encoded.
%%
%This is particularly powerful since we can, almost immediately detect, see if there are one, two, three of four objects in a group - an ability that we possess from early on in our development \cite{dehaene2011number,ware2019information}.
%%
%We consider similar assessments harder to make when colors are used.
%%
%For categorical variables, semantic icons can be used to represent categories.
%%
%Here, ideally, the icon design should be such that the underlying category is immediately recognizable and pre-attentively \textit{distinct}.
%%
%The benefits for semantic icons are two-fold: (1) the simplicity of design allows for easy recognition, as well as comparison of and the estimation of difference in values between multiple regions thus leveraging \textit{pre-attentive} processing, and (2) relieve the reliance of colors so that the data encoded in colors are easier to interpret.

% \note{Maybe have an example}

However, we also highlight some trade-offs of employing semantic icons.
%
First, by mapping continuous variable to discrete values, we effectively trade accuracy for speed and ease. 
%
Subtle differences between values will be lost, since different data values from two regions might both be mapped to, say, a set of two water droplet icons.
%
Second, by the nature of the encoding where data is encoded in the number of icons, the data must have a meaningful value to which regions with zero icons can be mapped.
%
Since regions with two icons will be perceived to be ``twice as much'' in comparison to regions with one icon, continuous variables whose range, for example, do not cross zero, may suffer from this encoding.


\subsubsection{Confidence encoding to address the modifiable areal unit problem (MAUP)}
\label{subsec:maup}
The modifiable areal unit problem (MAUP) is a well-documented problem in cartography and spatial data analysis, and is shown to have serious implications for decision-making \cite{fotheringham1991modifiable, manley2021scale, nelson_evaluating_2017}.
%
Broadly speaking, MAUP is an issue in geospatial data analysis where the results of an analysis can be influenced by the shapes, sizes and orientations of the polygons used to represent the data.

For example, a study on crime rate may reach different conclusions on whether the data is based on neighborhood, city, or county boundaries.
%
As Butkiewicz et. al. describe, MAUP can be broadly separated into two components: scale effect and aggregation effect \cite{butkiewicz_alleviating_2010}.
%
The former, the scale effect, relates to the fact that different spatial scales can dilute and obscure finer details when aggregated to a larger scale.
%
The latter, the aggregation or zoning effect, relates to the fact that particular spatial discretizations of the underlying space can lead to different patterns in the data.
%
In other words, different observations can be obtained by delineating data into different spatial blocks, even if the underlying data is the same.
%
Gerrymandering is one such example.

As we consider spatial binning to combine data of varying spatial discretizations, MAUP becomes a concern in visualization as well. 
%
Given entirely different data, MAUP can lead to the same aggregated result  in a given HexTile.
% \begin{wrapfigure}{r}{0.7\linewidth}
    \begin{figure}[htb]
    \centering
    \includegraphics[width=0.9\linewidth]{MAUP2}
    \caption{A motivating example for when MAUP interferes with visualization. Suppose a boundary between two regions in the original data bisects a HexTile domain. Note that in both cases, since the area within each HexTile is similar, a spatially-weighted average would report the same value of 70, despite the inherent variability in the data. The weighted mean is an incomplete picture of the data.}
    \label{fig:MAUP}
    \end{figure}
% \end{wrapfigure}
%
\cref{fig:MAUP} illustrates this by way of example.
%%
%Notice here that due to separate spatial discretization that HexTiles or any other spatial binning would introduce, aggregation within the data results in the loss of information.
%
Past works in visualizations have addressed MAUP including work by Butkiewicz et al. \cite{butkiewicz_alleviating_2010}, but their work is a system to assist in geospatial data analysis in the face of MAUP, and not a geospatial visualization encoding design with explicit consideration of MAUP.
%
With HexTiles, we explicitly encode \textit{confidence} in each HexTile component to address MAUP.

Specifically, we calculate a spatially-weighted variance for each of the variables we encode for each HexTile region.
% 
Then, we express that via Value-Suppressing Uncertainty Palettes (VSUP) \cite{2018_uncertainty_palettes}, inner ring thickness, and opacity, depending on the variable (details in \cref{sec:workflow}).
%
The user can therefore interpret the variability of the data within each HexTile, so that variability information is not lost in the visualizations.
 

% \fix{MAUP is a problem in geospatial data analysis where the results of an analysis can be influenced by the spatial scale of the data. This is because the spatial scale can affect the spatial patterns of the data, which in turn can affect the results of the analysis. The MAUP can lead to biased results and incorrect conclusions, which can have serious implications for decision-making. Often, MAUP is considered in statistical analysis of differently spatially discretized data, but we argue that
% it is also a problem in geospatial visualization as well. To address this, we calculate a confidence value via spatially and/or density weighted average within each HexTile, which can be used to interpret the variability of the data within each HexTile. This information can be lost in geospatial visualizations, but with HexTiles, we aim to make this information explicit.}

\subsection{Implementation}
% \\
% \note{Need a visual example of a HexTile, much like Lucas' report does. }

 \begin{figure}
   \begin{center}
        \includegraphics[width=0.8\linewidth]{HexTileAnatomy}
   \end{center}
   \vspace{-0.15in}
   \caption{The anatomy of a single HexTile and components in which data can be encoded: (1) the HexTile color, (2) the HexTile inner ring, and (3) the overlaid semantic icons (the design will vary with application)}
   \label{fig:HexTilesAnatomy}
   \vspace{-0.15in}
 \end{figure}
%\begin{figure}[htb]
%% \begin{figure}[]
%  \centering
%  % the following command controls the width of the embedded PS file
%  % (relative to the width of the current column)
%  % replacing the above command with the one below will explicitly set
%  % the bounding box of the PS figure to the rectangle (xl,yl),(xh,yh).
%  % It will also prevent LaTeX from reading the PS file to determine
%  % the bounding box (i.e., it will speed up the compilation process)
%  % \includegraphics[width=.95\linewidth, bb=39 696 126 756]{sampleFig}
%  %
%  %\parbox[t]{.9\columnwidth}{\relax
%  %         For all figures please keep in mind that you \textbf{must not}
%  %         use images with transparent background! 
%  %         }
%  %%
%  \label{fig:HexTilesAnatomy}
%  \caption{The anatomy of HexTiles}
%\end{figure}

The HexTiles design (\cref{fig:HexTilesAnatomy}) is a visualization encoding design to encode three geospatial variables at once, including our weighted-variance defined confidence to address MAUP, via a carefully designed superimposition.
%
The main software we leverage in HexTiles are the hierarchical hexagon-based spatial binning system \texttt{H3} \cite{uberh3} and the \texttt{WebGL}-accelerated \texttt{deck.gl} \cite{wang_deckgl_2019}.

\texttt{H3} is a hierarchical hexagon-based spatial binning system developed by Uber \cite{uberh3}. 
%
Unlike rectangular spatial bins, hexagon bins cannot simply be represented by a row and column due to the nature of hexagonal tessellation which has three degrees of direction rather than two. 
%
\texttt{H3} simplifies the location encoding as well as size encoding of hexagons by assigning each hexagon to an ID based on its location and size (resolution).
%
\texttt{H3} also simplifies the process of mapping the basemap zoom level (we use \texttt{Mapbox} in our implementation) to the appropriate hexagon tile size, since \texttt{H3} defines discrete resolutions that approximately subdivides a hexagon into seven smaller hexagon in the next resolution in its hierarchy. 
%
We refer the readers to the \texttt{H3} documentation \cite{uberh3} for more details.
%
In HexTiles, we use \texttt{H3} to define the hexagonal regions, or HexRegions, in which each tile is placed in.
%
Finally, in further pursuit of \textbf{DC3}, we allow users to \textit{zoom} into regions of interest, to discern finer details in the data, decreasing the hexagon size at each level.
%
This both mitigates the scale effect for MAUP, and ensures that HexTiles adheres to the \textit{information seeking mantra} of visualization design \cite{shneiderman_eyes_1996}. 
% \textcolor{red}

% \section{Design and Workflow}\label{sec:design}
\section{HexTiles Workflow}\label{sec:workflow}
%% \\
%% \note{Need a visual example of a HexTile, much like Lucas' report does. }

% \begin{figure}
%   \begin{center}
%        \includegraphics[width=0.8\linewidth]{HexTileAnatomy}
%   \end{center}
%   \vspace{-0.15in}
%   \caption{The anatomy of a single HexTile and components in which data can be encoded: (1) the HexTile color, (2) the HexTile inner ring, and (3) the overlaid semantic icons (the design will vary with application)}
%   \label{fig:HexTilesAnatomy}
% \end{figure}
%%\begin{figure}[htb]
%%% \begin{figure}[]
%%  \centering
%%  % the following command controls the width of the embedded PS file
%%  % (relative to the width of the current column)
%%  % replacing the above command with the one below will explicitly set
%%  % the bounding box of the PS figure to the rectangle (xl,yl),(xh,yh).
%%  % It will also prevent LaTeX from reading the PS file to determine
%%  % the bounding box (i.e., it will speed up the compilation process)
%%  % \includegraphics[width=.95\linewidth, bb=39 696 126 756]{sampleFig}
%%  %
%%  %\parbox[t]{.9\columnwidth}{\relax
%%  %         For all figures please keep in mind that you \textbf{must not}
%%  %         use images with transparent background! 
%%  %         }
%%  %%
%%  \label{fig:HexTilesAnatomy}
%%  \caption{The anatomy of HexTiles}
%%\end{figure}

%The HexTiles design (\cref{fig:HexTilesAnatomy}) is a visualization encoding design to encode three geospatial variables at once, including our weighted-variance defined confidence to address MAUP, via a carefully designed superimposition.
%%
%The main software we leverage in HexTiles is the hierarchical hexagon-based spatial binning system \texttt{H3} \cite{uberh3} and the \texttt{WebGL}-accelerated \texttt{deck.gl} \cite{wang_deckgl_2019}.

%\texttt{H3} is a hierarchical hexagon-based spatial binning system developed by Uber \cite{uberh3}. Unlike rectangular spatial bins, hexagon bins cannot simply be represented by a row and column due to the nature of hexagonal tessellation which has three degrees of direction rather than two. \texttt{H3} simplifies the location encoding as well as size encoding of hexagons by assigning each hexagon to an ID based on its location and size (resolution).
%%
%\texttt{H3} also simplifies the process of mapping the basemap zoom level (we use \texttt{Mapbox} in our implementation) to the appropriate hexagon tile size, since \texttt{H3} defines discrete resolutions that approximately subdivides a hexagon into seven smaller hexagon in the next resolution in its hierarchy. 
%%
%We refer the readers to the \texttt{H3} documentation \cite{uberh3} for more details.
%%
%We use \texttt{H3} to define the hexagonal regions, which we call HexRegions, on which HexTiles are placed.

%% In programming library implementations of H3, the calculation of the ID is abstracted away from the programmer, allowing one to easily map a single hexagonal region (in what we describe as a HexRegion) to its data properties via the HexRegion’s ID. 

%For HexTiles, there are two main workflow steps which need to be carefully considered shown in \cref{fig:workflow}: \textbf{Data Aggregation} and \textbf{Data Encoding}.
%%
%Note that these steps are required for each HexRegion and for each zoom level that \texttt{H3} defines to adhere to the \textit{overview first, zoom and filter, then details-on-demand} principle of visualization design \cite{shneiderman_eyes_1996}.
%% In designing HexTiles, the hex encoding function is a key design decision.
%%
%This function represents the mapping between the set of data found within each HexRegion and how we encode that data using the various channels of our HexTile visual encoding, like base color, inner ring, and the semantic icons.
%%
%The choice of the appropriate encoding funcion is crucial, and we outline guidelines for this in this section.
% In a simple case, $f$ can be a colormap mapping some continuous varialble to a range of colors
% $f:$ HexRegion $\to\{\mathbb{R},\mathbb{Z}\}$ 
% maps HexRegions and a particular variable to set of real numbers (in the case of color) or an integer (in the case of icons)

% , in which a HexRegion is a hexagonal-shaped bin that contains a set of datapoints defined within the bounds of the bin; the datapoints themselves may have multiple properties (such as name, height, and sex). The range of the hex encoding function is the set of integers and real numbers that define a characteristic of the HexTile, such as color, height, thickness, or amount of icons. 

For HexTiles, there are two main workflow steps which need to be carefully considered shown in \cref{fig:workflow}, \textbf{Data Aggregation} and \textbf{Data Encoding}, to transform raw data into a HexTiles visualization.
%
Note that these steps are required for each HexRegion and for each zoom level that the designer defines.

\begin{figure*}[t]
  \centering
  % the following command controls the width of the embedded PS file
  % (relative to the width of the current column)
  \includegraphics[width=0.85\linewidth]{Hex_workflow}
  % replacing the above command with the one below will explicitly set
  % the bounding box of the PS figure to the rectangle (xl,yl),(xh,yh).
  % It will also prevent LaTeX from reading the PS file to determine
  % the bounding box (i.e., it will speed up the compilation process)
  % \includegraphics[width=.95\linewidth, bb=39 696 126 756]{sampleFig}
  %
  %\parbox[t]{.9\columnwidth}{\relax
  %         For all figures please keep in mind that you \textbf{must not}
  %         use images with transparent background! 
  %         }
  %%
  \caption{HexTiles encoding workflow}
  \label{fig:workflow}
\end{figure*}

% Broadly, this function operates in two steps: data aggregation and data mapping.
\subsection{Data Aggregation}
%
\textbf{Data aggregation} describes the process of combining, or aggregating, the set of data within a HexRegion into one \textit{aggregate data point} that can be represented by a single HexTile. 
%%
%Specifically, a data point is a collection of attributes associated to a position in geographical space. 
%
This aggregation can be defined in a number of ways, including but not limited to a simple mean, spatially-weighted mean (i.e. a mean that takes into account the amount of overlap between the HexRegion and the data to weigh those data which are more present in the HexRegion), or any other function depending on the particular use case. 
%
The spatially-weighted mean is particularly useful for data formats that have attributes associated to polygonal features instead of a single position (such as GeoJSON data) since a HexRegion can span multiple polygons. 
%
In such cases, the data values of the polygons that the cell covers are spatially interpolated based on the area of its containing polygon that intersects with the hexagonal cell. 
%
It is worth noting that in cases where attributes vary in density between polygons, such as the percentage of the population who are Democrat where urban and rural areas vary widely in population, data values are also interpolated based on the density of the data collection in addition to its spatial weight. 
%
% \note{A simple graphic would be good here. I think we can show a HexRegion that includes polygon boundaries and show the aggregation. You also introduce some terminology like `property`, `datapoint`, `aggregated datapoint`. So a workflow type image here would help} 
%
% This aggregation is done through an aggregation function, the domain of which is a set of data points within a HexRegion and the range of which is a single data point that may or may not have the same attributes as the original set of datapoints, depending on what the user deems relevant or aggregable. 
%

Lastly, as discussed in \cref{subsec:maup}, we calculate a \textit{confidence} value for the \textit{aggregated data point}.
%
In other words, this confidence must address the question, ``How faithful is the aggregated data point to the actual data within the HexRegion?''
%
We adopt a weighted variance calculation; however, metrics like Moran's I that are used in spatial data analysis\cite{parenteau_modifiable_2011} can be used as well.
%
%Additionally, this aggregation function would ideally be done on multiple resolutions. 
%%
%In other words, on one iteration an aggregation function would be done on all HexRegions of one size, and on another iteration, the sizes of HexRegions would be changed to a smaller or larger size and the aggregation function done again on all HexRegions of this new size. 
%%
%This multiresolution aggregation would allow multiple levels of details during user interaction, for example by enabling more precise information to be revealed on zooming in on a screen through transitioning to smaller HexRegions.
%

%\begin{figure}[htb]
%  \centering
%  % the following command controls the width of the embedded PS file
%  % (relative to the width of the current column)
%  \includegraphics[width=\linewidth]{raster}
%  % replacing the above command with the one below will explicitly set
%  % the bounding box of the PS figure to the rectangle (xl,yl),(xh,yh).
%  % It will also prevent LaTeX from reading the PS file to determine
%  % the bounding box (i.e., it will speed up the compilation process)
%  % \includegraphics[width=.95\linewidth, bb=39 696 126 756]{sampleFig}
%  %
%  %\parbox[t]{.9\columnwidth}{\relax
%  %         For all figures please keep in mind that you \textbf{must not}
%  %         use images with transparent background! 
%  %         }
%  %%
%  \caption{Rasterizing polygons}
%  \label{fig:raster}
%\end{figure}

\subsection{Data Encoding}
\textbf{Data Encoding} describes the step where the set of \textit{aggregate data points} is mapped to the set of three HexTile components.
% We may select one attribute to map to HexTile height and another attribute to the number of icons. The collection of a single HexTile characteristic then becomes a channel in which the user can analyze information about one attribute of the dataset. For example, the height channel of HexTiles will represent one attribute of the dataset and the color channel will represent another attribute.
The choice of the specific mapping from the set of \textit{aggregate data points} to the HexTile components is crucial in providing an intuitive visualization.
%
%Multiple considerations must be made on the macrolevel and microlevel of visualization design.
%%
%On the macrolevel, general encoding choices must be ensured as to not cause too much visual clutter or ambiguity, as with the case of color schemes that are too similar or encoding too much information via the same channel type. 
%%
%We argue that using diversified channel types is more advantageous than encoding multiple parameters using one channel type even if they share similar traits. 
%%
%For example, instead of encoding two continuous attributes in two channels of color, it would be easier for a general user to interpret the two parameters if they were encoded on separate channel types, e.g. one using color and another using thickness.
%%
%On a microlevel, appropriate channel type and scaling must be properly chosen to ensure as minimal cognitive load on the user as possible and leverage pre-attentive processing. In some cases, there may be a tradeoff between the macro- and microlevel design considerations: for example, two parameter types may best be encoded in color, but one may have to be encoded in an alternate channel like icon amount to avoid using the same channel.

In \cref{table:encoding}, we present guidelines for what data attributes are ideal for each HexTile characteristic. 
%
We suggest that the base HexTile color be encoded to the most important attribute that one would like to guide the user's attention to. 
%
This is due to the fact that the base color takes up the most visual real estate and thus attracts the most attention from the user.
%
Icons should encode to the data attribute whose spatial distribution is of interest and the variable whose interaction with the other variables is of interest.
%
The speed at which the user can readily interpret the number of icons is the main advantage of this channel.
%
Here, the icon object itself should be immediately recognizable and unambiguous so that users can easily interpret what they represent. 
%
From our experiments, icons should not be used for variables if the resulting visualizations have many regions with many icons, as this can introduce visual clutter due to its 3D nature.
%
Lastly, the inner ring color should be used for secondary variables that are important for contextualizing the other attributes.
%
We also note that that some experimentation with colormaps for the base color and the inner ring will be required to ensure that each channel is best discriminated by the user.
%
If poorly chosen, the user may have difficulty interpreting the visualization, or worse, may misinterpret the data.

%Ring color and ring thickness can encode to secondary attributes that are important for contextualizing spatial distributions of the primary and distribution/categorical attributes. For ring thickness, we suggest choosing a data attribute that would justify an increase of visual real estate, such as confidence or error, since an increase of visual real estate would draw more user attention to regions with thicker rings. Height can be utilized for a fifth data attribute, but should encode to elevation or a similar data attribute that minimizes cognitive translation. In general however, height introduces occlusion and thus should be used minimally.


\newcolumntype{Y}{>{\raggedright\arraybackslash}X}
\begin{table*}[t]
\centering
\begin{tabularx}{\linewidth}{|c|c|Y|Y|}
\hline
\multicolumn{2}{|l|}{\textbf{HexTile Encoding}} & \textbf{Variance Encoding} & \textbf{Appropriate Attributes} \\ \hline
\multicolumn{2}{|l|}{Base Color} & Value-Suppressing Uncertainty Palette\cite{correll2018value}. We treat \textit{confidence} as \textit{uncertainty} in the language of the original work. Higher \textit{confidence} implies lower \textit{uncertainty}  & Primary attribute to guide user attention \\ \hline
 \multicolumn{2}{|l|}{Semantic Icons} & Opacity - more opaque for higher confidence so that these are more visually prominent  & Attribute where the spatial distribution or interactions with the other variables is the focus. Must have a reasonable value to which zero can mapped. \\ \hline
 \multicolumn{2}{|l|}{Inner Ring} & Thickness - thicker ring for higher confidence so that these are more visually prominent  & Secondary variables \\ \hline
 % Inner Design & Ring Thickness & Secondary attribute --- in particular when the assumption that some values ought to take up more visual real estate holds like uncertainty \\ \hline
 % \multicolumn{2}{|c|}{HexTile Height} & Elevation is a sensible option, but can encode any attribute. However should be used with caution as it can be distract from other variables and cause occlusion (further discussed in \cref{sec:feedback}). \\ \hline
\end{tabularx}
\caption{Guideline for choosing the mapping between data variables and HexTile components described in \cref{fig:HexTilesAnatomy} }
\label{table:encoding}
\vspace{-0.2in}
\end{table*}


\note{ So for hex, we have height, hex color, [ring width, inner ring color], icons, to encode data. Water we use inner ring color (mostly because the values are similar across so it doesn't introduce visual clutter), and for wildfire confidence we use ring width because this is a variable that we think that some values ought to have less visual impact than others. (Similar to \cite{correll2018value}). Finally we use icons for variables where we are interested in spatial
distribution and want to speed up / east the process of comparing different regions and its values. This overall structure for HexTiles should come through here, but we shouldn't reference the variables that are specific to each dataset like `power` or `demand units`, though I think elevation is fine. I'll send you my mental map for the designs on discord.}

\section{Case Studies}\label{sec:cases}
To evaluate HexTiles, we consider two data sets: (1) 2020 Presidential Election and demographics data in Texas, and (2) Outputs from CalSim3, the California Water Resources Simulation model of the Central Valley Project (CVP) and State Water Project (SWP) \cite{draper_calsim_2004}.
%
In particular, we compare HexTiles to Square-glyphs, another superimposition-based multivariate visualization technique with past use in geospatial data visualization and Cartography as discussed in \cref{sec:relworks}.

\subsection{Election data}
We consider the 2020 Presidential Election and demographics data in Texas, which includes the number of votes for each candidate in each precinct\cite{election_data}, as well as demographic data for each county subdivision\cite{census_2020}.
%
In particular, we focus on three variables: (1) Population density per squared kilometer, (2) Percent people of color, and (3) Percent democrat lead.
%
This use case is described in~\cref{subsec:election_usecase}.
%%
%As outlined in \cref{sec:relworks}, Square-glyphs, encode data in the size of the squares. 

\subsection{CalSim3 data}
Given the large set of data that CalSim3 can produce, with consultation with domain experts, we take a timestep from its output data for one simulation describing one alternate scenario using the same operational rules as baselines but with modified hydrology inputs reflecting a possible future climate condition. 
%

We consider three variables: (1) Unmet water demand, (2) Difference of unmet water demand with respect to historical baselines, and (3) Groundwater levels.
%
Unmet water demand is calculated at each \textit{demand unit}, a spatial discretization that CalSim3 uses to represent water use in the Sacramento and San Joaquin Valleys\cite{CalSim3}, where the difference between \textit{net water delivery} and \textit{water demand} is taken.
%
Furthermore, we also consider the difference in unmet water demand with respect to historical baselines as another important metric, as noted by our domain experts.
%
Lastly, we also take into consideration groundwater levels, as it plays an important role in the management of water resources in California's Central Valley.
%
Described using a different spatial discretization, groundwater levels report the amount of groundwater present in each location.
%
This use case is described in~\cref{subsec:water_usecase}.

\begin{figure*}[h!]
    \centering
    \includegraphics[width=\textwidth]{ElectionUseCase}
    \caption{2020 Presidential Election and demographics data in Texas. (1) Population density, (2) Percent people of color, (3) Percent democratic lead, (4) Square-glyph visualization of the three variables, (5) HexTiles visualization of the three variables without variance encoding, and (6) HexTiles visualization of the three variables with variance encoding.}
    \label{fig:usecase_election}
\end{figure*}
\begin{figure*}[h!]
    \includegraphics[width=\textwidth]{ElectionComparison}
    \caption{The 2020 Presidential Election and demographics data over Houston, TX encoded in (1) Square-glyphs, (2) HexTiles without variance encoding, and (3) HexTiles with variance encoding}
    \label{fig:usecase_election_comparison}
    \vspace{-0.15in}
\end{figure*}

\begin{figure*}[h!]
    \includegraphics[width=\linewidth]{MAUPUseCase}
    \caption{The same region of the election data use case near Dallas, TX encoded in HexTiles with (1) HexTiles without variance encoding, (2) HexTiles with variance encoding, and (3) HexTiles with variance encoding at a higher zoom level. Notice that the hexagonal region approximately subdivides into seven smaller hexagonal regions.}
    \label{fig:usecase_maup}
\end{figure*}

\begin{figure*}[h!]
    \centering
    \includegraphics[width=\textwidth]{WaterUseCase}
\caption{CalSim3~\cite{CalSim3} output for a particular water allocation scenario at a particular timestep. (1) Unmet water demand at each demand unit, (2) The difference of unmet water demand with respect to historical baselines, (3) groundwater levels, (4) Square-glyph visualization of the three variables, (5) HexTiles visualization of the three variables without variance encoding, and (6) HexTiles visualization of the three variables with variance encoding.}
    \label{fig:usecase_water}
    \vspace{-0.15in}
\end{figure*}


\subsection{Visualizations}
Here, we describe the resulting visualizations and the specifics for the two use cases, as well as the impact of the variance encoding for HexTiles, in order.

\subsubsection{2020 Presidential Election Data Visualization in Texas}
\label{subsec:election_usecase}
Using the 2020 Presidential Election and demographics data in Texas, we present the resulting visualizations for Square-glyphs, HexTiles without variance encoding, and HexTiles with variance encoding in~\cref{fig:usecase_election}.
%
For a detailed comparison, we also show the visualizations over the same area near Houston, TX in~\cref{fig:usecase_election_comparison}.
% 
For Square-glyphs, data is encoded in the size of the squares as introduced prior. 
%
However, for diverging data like the percent Democrat lead, in the interest of staying genuine to its original design, we divide the data into two separate squares, one square representing the lead for each party.
%
If combined into a single square, the size of the square could be misleading.
%
Thus, from the top-left square clockwise, the squares and its size encode, (1) Percent Democrat lead, (2) Percent Republican lead, (3) Population density, and (4) Percent people of color.
%
For HexTiles, its base color encodes percent Democrat lead, its inner ring the percent people of color, and the semantic icon in the center the population density.
%
The according legends are shown on the right.

\subsubsection{Water Distribution Visualization from CalSim3}
\label{subsec:water_usecase}
Using the CalSim3 data, we present the resulting visualizations for Square-glyphs, HexTiles without variance encoding, and HexTiles with variance encoding in~\cref{fig:usecase_water}.
%
For the Square-glyphs, similar to the percent Democrat lead in the election data, we divide the difference with respect to historical baselines into two separate squares, one for positive differences and one for negative differences.
%
From the top-left square clockwise, the squares and its size encode, (1) Scenario unmet water demand (i.e. unmet water demand for this specific scenario), (2) Groundwater levels, (3) Positive difference of unmet water demand with respect to historical baselines, and (4) Negative difference of unmet water demand with respect to historical baselines.
%
For HexTiles, its base color encodes groundwater levels, its inner ring the difference of unmet water demand with respect to historical baselines, and the semantic icon in the center the scenario unmet water demand.
%
The according legends are shown on the right.


\subsubsection{Variance encoding to alleviate the modifiable areal unit problem}
To highlight the impact of the variance encoding and the more faithful representation of the data it enables, we describe a specific example using the election data in \cref{fig:usecase_maup}.
%
The region of interest here is near Dallas TX. 
%
With the variance encoding, initially the user is shown a HexTile with unsaturated blue-purple base color (somewhat Democrat) and a dark green (high percent people of color) but thin inner ring.
%
The unsaturated base color, as well as the thin inner ring communicates, that the spatially-weighted mean alone is not a \textit{faithful} representation of the data.
%
Zooming into the region, the user sees the HexTiles subdivide seven into smaller HexTiles, each with a different base color and inner ring thickness.
%
In fact, in this region near Dallas, there is a diverse mix of areas with varying levels of Democrat lead and percent people of color.
%
The differences in characteristics of urban and rural areas near Dallas and Forth Worth are more \textit{faithfully} communicated with the variance encoding.
%
Note that these insights are entirely absent without the variance encoding.
%


\section{Evaluation}\label{sec:feedback}
\begin{figure*}[h!]
    \centering
    \begin{subfigure}[b]{0.24\linewidth}
        \centering
        \includegraphics[width=\textwidth]{T1Error}
        \caption{Error for Task 1}
        \label{subfig:T1Err}
    \end{subfigure}
    \begin{subfigure}[b]{0.24\linewidth}
        \centering
        \includegraphics[width=\textwidth]{T1TTI}
        \caption{Time taken for Task 1}
        \label{subfig:T1TTI}
    \end{subfigure}
    \begin{subfigure}[b]{0.24\linewidth}
        \centering
        \includegraphics[width=\textwidth]{T2Rate}
        \caption{Success rate for Task 2}
        \label{subfig:T2Err}
    \end{subfigure}
    \begin{subfigure}[b]{0.24\linewidth}
        \centering
        \includegraphics[width=\textwidth]{T2TTI}
        \caption{Time taken for Task 2}
        \label{subfig:T2TTI}
    \end{subfigure}
    \caption{Task performance metrics}
    \label{fig:tasks}
\end{figure*}

\begin{figure*}[h!]
    \centering
    \begin{subfigure}[b]{0.48\linewidth}
        \centering
        \includegraphics[width=\textwidth]{T1TLX}
        % \caption{Task }
        \label{subfig:T1TLX}
    \end{subfigure}
    \begin{subfigure}[b]{0.48\linewidth}
        \centering
        \includegraphics[width=\textwidth]{T2TLX}
        % \caption{T2TLX}
        \label{subfig:T2TLX}
    \end{subfigure}
    \caption{Hart and Staveland's NASA-TLX scores for Task 1 and Task 2. Asked were Mental effort, Mental demand, Frustration, and Performance/accuracy of HexTiles with respect to Square-glyphs.}
    \label{fig:tasks_tlx}
    \vspace{-0.15in}
\end{figure*}

To evaluate the efficacy of our HexTiles design, we conducted a user study with 7 participants with two tasks, measuring accuracy and time-to-insight, and an expert review with three domain experts (DEs). 

\subsection{User Study}
We gathered seven participants, \textbf{P1}-\textbf{P7}: 2 male and 5 female; 5 PhD students, 1 MS student, and 1 BS student; with varying levels of exposure to visualization.
%
The user study was conducted with visualizations using the 2020 Presidential Election and demographics data as seen in \cref{fig:usecase_election}.
%%
%The three variables encoded were: (1) Population density, (2) Percent people of color, and (3) Percent democratic lead.

As the baseline, we used the existing superimposition-based method, Square-glyphs, and compared its performance against HexTiles.
%
For the sake of comparison, HexTiles without the variance encoding was used in the study.
%
In addition, to ensure fairness of comparison, we allowed participants to freely change elevation angles for both visulizations to account for the 3D nature of HexTiles.


We designed two tasks as follows: 
\begin{tightItemize}
    \item \textbf{Task 1:} Given a Square-glyph or HexTile, \textit{estimate} the value of the encoded variable.
    \item \textbf{Task 2:} Given a desired pair of variable ranges, \textit{find} a Square-glyph or HexTile that satisfies the condition. We use this task as a proxy to measure how well users can interpret and reason about the interactions between the encoded variables.
\end{tightItemize}
As an example, for Task 2, participants were asked to find and click on a tile with \textit{population density <25 per square km \textbf{and} percent people of color >50\%}.
%
For both tasks, participants were measured on accuracy and time-to-insight (TTI).
%
To define accuracy, for Task 1, we used the error between the estimated and actual value expressed as a percentage from the minimum and maximum of that variable, and for Task 2, we used the success rate.
%
To define TTI, we measured the time taken to complete each task.
%
The regions asked and the ranges asked were randomized for each participant, and 15 questions were asked for each task and for each visualization type, resulting in $15 \times 2 \times 2 = 60$ questions in total for every participant.
%
The four hypotheses we tested at a significance level of $\alpha = 0.05$ were:
\begin{tightItemize}
    \item \textbf{H1:} HexTiles will have a \textit{lower} error rate compared to Square-glyphs on Task 1.
    \item \textbf{H2:} HexTiles will have a \textit{shorter} TTI compared to Square-glyphs on Task 1.
    \item \textbf{H3:} HexTiles will have a \textit{higher} success rate compared to Square-glyphs on Task 2.
    \item \textbf{H4:} HexTiles will have a \textit{shorter} TTI compared to Square-glyphs on Task 2.
\end{tightItemize}
\subsubsection{Quantitative Results}
Quantitative results are shown in \cref{fig:tasks}.
%
We ran a Mann-Whitney U test for each hypothesis (\textbf{H1}, \textbf{H2}, \textbf{H4}) except for \textbf{H3} where we ran a two-sample Z test for proportions.
%
For \textbf{H1}, \textbf{H2}, and \textbf{H4}, we opted for the Mann-Whitney U test due to the non-normal distribution of the data.
%
At a significance level of $\alpha = 0.05$, we failed to reject the null hypothesis for all hypotheses.
%
However, at a p-value of 0.117 for \textbf{H3}, we note that this still points to a trend towards HexTiles having a higher success rate compared to Square-glyphs though not statistically significant.
%
Furthermore, for \textbf{H1}, even though we failed to reject the null hypothesis, we find a lower mean error for HexTiles compared to Square-glyphs.

One notable observation was the variances in error and TTI for Task 1, and TTI for Task 2, which were \textit{lower} for HexTiles compared to Square-glyphs.
%
Administering the Levene's test for equality of variances, we found that the variances were \textit{significantly different} for error in Task 1 and TTI in Task 2 at a significance level of $\alpha = 0.05$.
%
We thus conclude that HexTiles has a \textbf{more consistent performance and mental load for interpretation} compared to Square-glyphs.

\subsubsection{Qualitative Results}
To gather qualitative feedback, we asked participants to (1) describe their experience with each visualization type, (2) answer the NASA-TLX questionnaire to measure the mental load for each task and visualization type, and (3) evaluate the variance encoding in HexTiles to see if they gained additional insights that were not present without the encoding.
%
Results for the NASA-TLX survey are shown in \cref{fig:tasks_tlx}.

On all four categories of the NASA-TLX survey tested\textcolor{red}{\textbf{---}}mental demand, performance, effort, and frustration\textcolor{red}{\textbf{---}}all participants \textbf{P1}-\textbf{P7} reported significantly favorable scores for HexTiles compared to Square-glyphs.
%
Here, we also report a number of key recurring points that the participants reported during the qualitative feedback session on the usability of HexTiles compared to Square-glyphs.

\paragraph{``Color is easier to interpret than size''}
A main distinction between HexTiles and Square-glyphs is the use of color vs. size.
%
Where the HexTiles design uses the color of the tile's base or ring to encode the value of a variable, the Square-glyphs design uses the size of the square.
%
% Where in HexTiles base color and inner ring, color is used to encode the value of variable, in Square-glyphs, the size of the square is used. \textcolor{red}{\textbf{possibly reword: ""}}
%
A number of participants noted that the size of the square was harder to interpret compared to the color of the HexTiles, especially when the size of the square was small.
%
\textbf{P1} noted, \qu{when you zoom out on the square glyphs, the little squares disappear. You can't see them anymore, [\dots] and then you can't really compare the sizes of the squares.}
%
They further noted, ``with HexTiles, the sizes didn't really matter.''
% Similarly, \textbf{P3} noted, ``[estimating values with HexTiles] was easier because I don\'t have to compute what the size of square should look like. [\dots] for both Task 1 and Task 2, I don't need to do this translation.''
%
\textbf{P2} agreed as well: \qu{[HexTiles] definitely took less effort to interpret because I can just look the scales, look at the colors. [With Square-glyphs] you need to look at it and then kind of estimate what it is.}
%
\textbf{P5} also left a similar comment that they had to check the sizes for Square-glyphs, whereas for HexTiles, they could just look at the color, which made it substantially less demanding to interpret and locate specific tiles.
%
Another notable insight from \textbf{P6} is that due to the nature of a map and the perspective transform, Square-glyphs and their relative square sizes will depend on the distance to zoom level, which can be a source of confusion and imprecision.
%
The reported comparison between color and size was an interesting observation, since literature in visualization\cite{Munzner_2015} suggests that size is a more effective channel than color for ordered attributes.
%
However, we report that in the context of superimposition-based multivariate geospatial visualizations, size may be limiting if many such tiles are to be used, especially for interactive visualizations when the perspective transform can interfere with the size perception.

\paragraph{``Icons are mentally easy to interpret''}
Participants also reaffirmed the efficacy of the semantic icons in HexTiles.
%
With regards to the difference between HexTiles and Square-glyphs, \textbf{P4} said, \qu{The biggest difference between [the two] would be reading and finding the icons for the population. [\dots] Population density [in the icons], was the easiest to interpret.}
%
\textbf{P2} noted, \qu{The population density [encoding for HexTiles] was very intuitive.}
%
\textbf{P7} reaffirmed the semantic icons, \qu{for population, definitely, HexTiles are better, just because you have that 3D icon, so it was really easy to tell.}
%
We found evidence that the semantic icons were useful in guiding user attention to specific variables over others.

\paragraph{``HexTiles was easier for understanding multiple variables at once''}
% {\textbf{possibly reword: "easier for understanding multiple variables at once"}}''}
Participants also noted that HexTiles made it easier to interpret multiple variables at once and to ascertain relationships between them.
%
\textbf{P7} said, \qu{I feel like HexTiles are really good at combining, [for example] when you're combining the political outcomes as well as the population density.}
%
\textbf{P5} remarked that Square-glyphs might be better if one focuses on one variable, but \qu{if I want to handle the three types together, HexTiles are better.}

\paragraph{``HexTiles provides a more visually continuous visualization that's easier to interpret and to locate a particular tile''}
Especially for Task 2, when participants were asked to find a tile of a given characteristic, participants noted that the mental demand was lower for HexTiles mostly due to the more continuous nature of the visualization design.
%
\textbf{P1} said, \qu{[Square-glyphs] all look alike when you zoom out. So you can't really see a big picture of [\dots] all the tiles, which is important when you're looking for a tile.}
%
Since in the Square-glyphs, the four squares of which the glyph is composed are separated, users mentioned the increased mental load to locate particular tiles.

\paragraph{``The variance encoding in HexTiles is useful''}
Participants were asked to evaluate the variance encoding in HexTiles and to see if they gained any additional insights that were not present without.
%
Many pointed to the variance encoding as a useful feature especially for exploring the data over urban areas like Dallas and Houston, as well as regions between urban areas, since users were able to better understand the nuanced transition of political lead from Democratic to Republican, and vice versa.
%
\textbf{P1} said, \qu{you have these places that are sort of in the middle between two regions [\dots] the variance is much higher within them.}
%
Participants generally agreed that the variance encoding was useful for adjusting confidence in the aggregated data presented as HexTiles.




% \paragraph{``HexTiles base color and inner ring color can compete which makes the inner ring hard to interpret''}

%\paragraph{``Especially with the variance encoding, HexTiles can be visually complex and overwhelming''}
%Some participants noted that HexTiles, especially with the variance encoding, can be visually complex and overwhelming.
%%
%\textbf{P5} said even though with time, one can get used to it and that it became better, the color is a lot at the start to take in.
%%
%\textbf{P1} put it nicely, ``[Square-glyphs] were easier on the eyes to look at because it was more simple. But HexTiles, [with the variance encoding] was more informative.''


\subsection{Expert Review}
We gathered feedback from three domain experts who work in water resource management all with expertise in ecology and hydrology and regularly use geospatial visualizations in their line of work.
%
% We denote these three domain experts as \textbf{DE1}-\textbf{DE3}.
%
For all \textbf{DE1}-\textbf{DE3}, we presented the visualizations in~\cref{fig:usecase_water}.
%
In these DE interviews, we focused on (1) the usability and clarity of HexTiles in the context of communicating water management data, (2) the efficacy of the variance encoding to alleviate MAUP, and (3) the limitations.
%
%In addition to DE interviews, we also gathered feedback via informal interviews from two PhD students in the field of visualization.
%%
%We will mainly focus on the DE interviews for usability feedback, but we also include results from our informal interviews.

% \vspace{-0.15in}

\subsubsection{Results}
All three DEs responded positively on the efficacy of communicating multivariate data via HexTiles. 
%
\textbf{DE1} described our design as \qu{visually fascinating and intriguing}, \textbf{DE2} described that the design \qu{definitely works}, and \textbf{DE3} remarked that HexTiles \qu{certainly goes beyond anything [they've] come across before.}
%
The strengths of HexTiles, as expected from our initial design stage, lied in showing general patterns of multivariate data---however, of course, with the cost of abstracting out finer details as mentioned in \cref{sec:methods}.
%
\textbf{DE1} suggested that HexTiles can be combined with other metrics and visualizations to communicate finer details.
%
\textbf{DE3} said \qu{[HexTiles] is really good for conveying the patterns over the whole region,} while also noting that the design can \qu{generalize over} a region if the specific area of interest is
relatively fine.
%
They concluded by highlighting the value in \qu{ignoring some specific information, [\dots] in favor of this generalizable approach.}
%
\textbf{DE2} also supported HexTiles for other similar uses in hydrology to communicate multivariate data: \qu{I'm actually really intrigued about applying it to a different project we're working on [\dots] the Sacramento River.} 
%
With regards to the variance encoding, \textbf{DE2} noted \qu{I think it's really interesting [that] it allows you to look at valley-wide patterns, but warns you about interpreting certain results at that scale, and [that] you need to zoom in more.[\dots] it's incredibly powerful and a ton of data.}
%
While the majority of comments were positive,  these interviews, in addition to those outlined during the user study, also reiterated a number of limitations with HexTiles which we list next.

\subsection{Limitations and Considerations}
We summarize the limitations and shortcomings of HexTiles as reported by our user study participants and domain experts.

%\textbf{Height:} The use of height (i.e. 2.5D) in spatial visualization and its efficacy is a recurring topic in visualization.
%%
%In our case, for example in \cref{fig:usecase_water}, we encode the land surface elevation as the height of the HexTiles.
%%
%A number of DEs, as well as PhD students during our informal feedback sessions noted that the height was one of the first information that ``popped out`` and that it distracted from interpreting other information.
%%
%While \textbf{DE2} mentioned that it may assist in their particular case since land surface elevation relates to the gravitational flow of water, all respondents agreed that it increased the complexity of the visualization, including the occlusion that it can cause.
%%
%As a general rule of thumb, as we discuss in \cref{sec:design}, the use of height in HexTiles should be done carefully, considering occlusion as well as the interpretability problems it may lead in 2.5D/3D visualizations.
%% 
%Although we include HexTiles height as a possible channel in which to encode data, especially given the feedback during the review process, we caution against it as a general rule of thumb.
%%
%Thus, we leave the detailed study of effective uses of HexTiles height as future work.

\noindent
\textbf{Color interactions: } When asked on the limitation of HexTiles, all participants in the user study noted that the base color and inner ring color can be difficult to interpret together, especially when their saturation is similar. 
%
For the election use case tested, the base color saturation at the middle of the colormap and the inner ring saturation at the end of the colormap were similar, causing difficulty in interpretation.
%
Since Square-glyphs do not inherently blend the encodings for the three variables, the visual separability was higher than the base color and inner ring color in HexTiles.
%
However, we reiterate that putting emphasis on certain variables over others is part of the design of HexTiles, and we see this as a consequence of that decision.
\\
\textbf{Visual complexity: } Some participants and domain experts noted that HexTiles, especially with the variance encoding, can be visually complex and overwhelming.
%
\textbf{P5} said even though they got used to the design and encodings with time, the color was a lot to take in, especially at the start.
% said even though with time they got used to the design and encodings, the color was a lot, especially at the start, to take in. \textcolor{red}{\textbf{possibly reword: s}}
%
\textbf{P1} put it nicely, \qu{[Square-glyphs] were easier on the eyes to look at because it was simpler. But HexTiles, [with the variance encoding] was more informative.}
%
\textbf{DE2} also remarked that the design can be overwhelming for some users who are not used to such complex visualizations.
%
The amount of information communicated and the visual complexity is a fine balance, and we acknowledge this limitation.

\noindent
\textbf{Levels of abstractions: }
Though \textbf{DE2} noted that this is not necessarily a limitation of HexTiles, but a \qu{limitation for any spatial data or almost any spatial data}, we nevertheless note that this is a limitation of a superimposition-based visualization.
%
While visualizing general behavior is a useful and crucial feature, as our domain experts (\textbf{DE2}, \textbf{DE3}) confirmed, we acknowledge that finer investigation may be better suited with a different visual design and encoding.
%
Adaptive task-based hexagonal resolution is another possible future direction for HexTiles.

\noindent
\textbf{Occlusion of underlying terrain: }
The HexTiles design drawn over a base map will occlude underlying natural and artificial features of the space.
%
While many glyph designs encounter this, for applications that explicitly require the underlying features like the terrain, etc. to be visible, HexTiles may not be the best choice.


\section{Conclusion}\label{sec:conclusion}
We introduced HexTiles with semantic icons, a domain-agnostic superimposition-based visualization design for multivariate geospatial data.
%
With its design, qualitative results show that HexTiles is able to more easily communicate multivariaate data with MAUP with in mind as seen by the NASA TLX survery results, albeit at small sample size.
%
Further our semi-structured interview results point to the efficacy of semantic icons along with its design also puts the visualization designer in the driver's seat for guiding user attention to specific variables.
%
Our quantitative results also showed some evidence that HexTiles allows for more accurate interaction assessment, and showed strong evidence that HexTiles is more consistent in its performance and mental load. 
%
We presented feedback from domain experts that echoed these findings.
%
We also list the limitations of HexTiles gathered through our evaluations -- though some are inherent to a superimposition-based design such as HexTiles.
%
Future work includes exploring the usability of HexTiles height for data encoding, adaptive task-based hexagonal resolution, consideration of temporal data (i.e., implementing HexTiles for spatiotemporal data and the necessary design considerations), and exploring the use of animation and further interaction to better support exploration and faithful sensemaking of multivariate geospatial data.

